\section{Quadratic Frobenius and mate lines}

Let $\cC\subseteq\PG(2,K)$ be a $\phi$-invariant $k$-arc.  Write
\[
  f=|\cC\cap\PG(2,F)|,
  \qquad k=f+2e,
\]
so that $\cC$ is the disjoint union of $f$ fixed points and $e$ nonfixed
conjugate pairs.

For every nonfixed point $P$, the line $P\phi(P)$ is fixed by $\phi$; we
call it the \emph{mate line} of the pair $\{P,\phi(P)\}$.  Conversely, every
nonfixed conjugate pair has a unique mate line.  A fixed line contains
$s^2+1$ points, of which $s+1$ are fixed, and therefore contains exactly
\[
  N=\frac{s^2-s}{2}
\]
nonfixed conjugate pairs.

\begin{definition}\label{def:legal-pair}
A \emph{legal conjugate-pair extension} of $\cC$ is an unordered pair
$Q=\{P,\phi(P)\}$ such that $P\ne\phi(P)$, $Q\cap\cC=\varnothing$, and
$\cC\cup Q$ is an arc.  Let $\Npair(\cC)$ be the number of such pairs.
\end{definition}

\begin{definition}[Alternate-orbit repair]\label{def:alternate-repair}
Let $A$ be an invariant arc, let $O\subseteq A$ be a selected nonfixed
conjugate orbit, and put $D=A\setminus O$.  An \emph{alternate-orbit repair}
of this deletion is a legal conjugate-pair extension $Q$ of $D$ with
$Q\ne O$.  The erased orbit $O$ is itself legal for $D$, so a lower bound
$\Npair(D)\ge r+1$ supplies at least $r$ alternate repairs.
\end{definition}

\begin{lemma}[Pair-extension test]\label{lem:pair-test}
Let $P\ne\phi(P)$ and suppose neither point belongs to $\cC$.  Then
$\cC\cup\{P,\phi(P)\}$ fails to be an arc if and only if at least one of the
following holds:
\begin{enumerate}[label=(\alph*),leftmargin=2em]
\item $P$ lies on a secant of $\cC$;
\item $\phi(P)$ lies on a secant of $\cC$;
\item the mate line $P\phi(P)$ contains a point of $\cC$.
\end{enumerate}
\end{lemma}

\begin{proof}
Every collinear triple in the enlarged set either contains one new point
and two old points, giving (a) or (b), or contains both new points and one
old point, giving (c).  The converse is immediate.
\end{proof}

The useful carriers are therefore the fixed lines disjoint from $\cC$.
Let $\cE$ be this set of empty fixed lines.

\begin{lemma}[Empty fixed lines]\label{lem:empty-lines}
The number of empty fixed lines is
\begin{equation}\label{eq:E}
 E=|\cE|=s^2+s+1-\left(f(s+1)-\binom f2+e\right).
\end{equation}
\end{lemma}

\begin{proof}
The $f$ fixed selected points occupy
$f(s+1)-\binom f2$ fixed lines.  Inclusion--exclusion is exact because no
three selected fixed points are collinear.  Each selected conjugate pair
occupies its mate line.  These $e$ lines are distinct and contain no fixed
selected point, since $\cC$ is an arc.  Subtract from the $s^2+s+1$ fixed
lines.
\end{proof}

Among the $\binom k2$ old secants, exactly
\[
 I=\binom f2+e
\]
are fixed: those joining two fixed selected points and the $e$ mate lines.
The remaining secants occur in two-element Frobenius orbits.  Their number
is
\begin{equation}\label{eq:M}
 M=\frac{\binom k2-I}{2}=fe+e(e-1).
\end{equation}
